\section{Conclusion}

This paper presented a mission planning simulation tool for the
in-space assembly of large spacecraft at Earth-Sun L4. The optimizer
formulates the campaign as a multi-objective dynamic programming
problem over a time-expanded network, with forward beam search
pruning to manage the exponential state space. Eight constraints
are enforced: assembly DAG precedence, EVA pair and IVA support
allocation, a 60\% crew duty cycle, scale-adaptive proximity
congestion, partial assembly work-in-progress tracking, pad
turnaround intervals, crew vehicle cost modeling, and cargo
vehicle proximity contribution during unloading.

Results across the scenarios presented demonstrate that the
choice of propulsion architecture (Chemical vs.\ NEP/Fusion)
has a larger effect on program cost and module mass than on
total mission duration, which is primarily determined by crew
rotation strategy and the number of concurrent cargo deliveries.
The Pareto frontier reveals that modest weight shifts toward
minimizing time require only small increases in launch count,
while cost-minimal solutions incur substantial duration penalties.

\textbf{Limitations.} The beam search is a heuristic and
does not guarantee global optimality. Delta-$v$ values are
computed from simplified Hohmann and WSB approximations rather
than full ephemeris-based trajectory optimization. Assembly
durations are treated as deterministic; no failure modes or
schedule margin is modeled.

\textbf{Future Work.} Higher-fidelity orbital mechanics
(e.g., patched-conic or GMAT integration), stochastic assembly
duration distributions, multi-site assembly staging through
cislunar Gateway, in-situ propellant production at L4, and
refueling depot optimization are natural extensions of this work.
